%% kapa: exam=1 round=33 year=2022 period=1 form=A subject=civil_law subject_ko=민법 seq=11 number=11
%% source: Q-Net 감정평가사 기출문제 (cst00309) artlSeq=5209997 file=raw/1008/5209997/2022년도 제33회 감정평가사 1차시험 1교시(민법,경제학,부동산학).pdf
%% conversion: kordoc 4.24.51 -> tools/gen1.py (auto); review=proofread-sonnet; flags=none
\begin{question}{11}
착오에 의한 의사표시에 관한 설명으로 옳지 않은 것은? (다툼이 있으면 판례에 따름)
\begin{choices}
  \item 대리인이 의사표시를 한 경우, 착오의 유무는 본인을 표준으로 판단하여야 한다.
  \item 착오가 표의자의 중대한 과실로 인한 때에는 표의자는 특별한 사정이 없는 한 그 의사표시를 취소할 수 없다.
  \item 착오로 인하여 표의자가 경제적인 불이익을 입지 않았다면 법률행위 내용의 중요부분의 착오라 할 수 없다.
  \item 상대방이 표의자의 진의에 동의한 경우, 표의자는 착오를 이유로 그 의사표시를 취소할 수 없다.
  \item 착오를 이유로 의사표시를 취소하는 자는 착오가 없었더라면 의사표시를 하지 않았을 것이라는 점을 증명하여야 한다.
\end{choices}
\end{question}
\answer{1}
